\addtocontents{toc}{\protect\setcounter{tocdepth}{1}}
\chapter[Definición del Proyecto]{Definición Proyecto}

\section[Objetivos del Proyecto]{Objetivos del proyecto}

\subsection{Objetivo general}

Desarrollar una solución móvil de apoyo integral a cuidadores de pacientes
oncológicos pediátricos para un centro hospitalario de alta complejidad, que
centralice herramientas de registro, seguimiento, consulta y orientación para
facilitar la continuidad del cuidado en el hogar.

\subsection{Objetivos específicos}

\begin{enumerate}
    \item Analizar las necesidades de información del cuidador principal y los
    flujos del cuidado domiciliario, a partir de los contenidos y criterios
    clínicos entregados por el equipo de Enfermería, para fundamentar la
    especificación funcional y no funcional del software.
    \item Diseñar la arquitectura de la solución móvil centrada en el cuidador,
    con su modelo de datos, servicios lógicos e interfaz de usuario, de modo que
    permita gestionar perfiles, registrar el seguimiento domiciliario y consultar
    el historial, priorizando la facilidad de uso y la seguridad de la
    información.
    \item Implementar los módulos que permiten al cuidador acceder a contenidos
    educativos validados, registrar los signos y síntomas del paciente, recibir
    orientación mediante alertas de semaforización clínica y gestionar
    recordatorios, con sincronización y respaldo en el servidor disponibles incluso
    sin conexión.
    \item Validar técnicamente el prototipo mediante pruebas automatizadas,
    verificando su conformidad con la especificación de requerimientos.
\end{enumerate}

\section[Ambiente de Ingeniería de Software]{Ambiente de Ingeniería de Software}

\subsection{Metodología de desarrollo}

El proyecto adopta Scrum de forma adaptada a un equipo de dos desarrolladoras y
a un plazo académico. Los requerimientos se expresan como historias de usuario
con criterios de aceptación, se estiman en puntos y se ordenan en un
\emph{product backlog} según su prioridad clínica y sus dependencias técnicas.
Scrum se elige porque permite entregar una versión funcional al cierre de cada
sprint, mostrarla a la contraparte y ajustar el alcance con lo aprendido.

El desarrollo se organiza en tres sprints. La Tabla~\ref{tab:sprints} resume su
foco y su calendario.

\begin{table}[H]
\centering
\small
\tablazebra
\caption{Organización del desarrollo en sprints.}
\label{tab:sprints}
\begin{tabularx}{\textwidth}{P{0.55}P{0.85}P{1.60}}
\toprule
\filaencabezado \encab{Sprint} & \encab{Periodo} & \encab{Foco} \\
\midrule
Sprint 1 & 26 ago -- 15 sep 2026 &
Flujo principal: acceso, cifrado de datos, gestión de pacientes, registro
clínico con clasificación de alertas e historial. \\
Sprint 2 & 16 sep -- 20 oct 2026 &
Rutina del cuidador: panel principal, exportación, biblioteca educativa,
asistente de orientación, preguntas frecuentes, recordatorios y listas de
verificación. \\
Sprint 3 & 21 oct -- 13 nov 2026 &
Integración: datos por paciente, aislamiento de la información, navegación
definitiva y funcionamiento sin conexión. \\
\bottomrule
\end{tabularx}
\end{table}

Cada sprint comienza con la selección de historias del backlog y termina con una
versión instalable (APK) que se revisa con la contraparte de Enfermería. La
contraparte valida la pertinencia de los contenidos, los textos y los criterios
de alerta, mientras que el docente guía supervisa el avance del proyecto.

\subsection{Técnicas y notaciones}

\begin{itemize}
    \item \emph{Historias de usuario} con criterios de aceptación, para expresar
    los requerimientos desde la perspectiva del cuidador.
    \item \emph{BPMN 2.0}, para modelar el proceso de seguimiento actual y el
    proceso futuro con la aplicación.
    \item \emph{UML 2.0}, en particular diagramas de casos de uso, para
    representar la interacción entre los actores y el sistema.
    \item \emph{Modelo entidad-relación}, para el modelo conceptual de datos, y
    esquema relacional de PostgreSQL para el diseño físico.
    \item \emph{Prototipos de alta fidelidad} en Figma, para validar la interfaz
    con la contraparte antes de construirla.
\end{itemize}

\subsection{Estándares de documentación}

La especificación de requerimientos se basa en IEEE~Std~830-1998
\citep{ieee_830_1998} y la documentación de pruebas en IEEE~Std~829-1998
\citep{ieee_829_1998}. Los atributos de calidad del producto toman como
referencia el modelo ISO/IEC~9126. En el código se sigue la guía de estilo de
Dart, verificada con flutter analyze, y los mensajes de
\emph{commit} siguen la convención \emph{Conventional Commits} en español.

\subsection{Tecnologías y herramientas de apoyo}

La aplicación se construye con Flutter \citep{flutter_2026}, que permite
compilar para Android e iOS desde una sola base de código en Dart y ofrece
componentes Material~3 listos para interfaces accesibles. Frente a Kotlin
nativo, evita duplicar el desarrollo para cada plataforma, y frente a React
Native, entrega un rendimiento más estable al no depender de puentes nativos.
Como backend se desarrolla una API REST propia con Node.js
\citep{nodejs_2026} y Express \citep{express_2026}, que persiste los datos en
una base relacional PostgreSQL \citep{postgresql_2026}. Ambos servicios se
empaquetan con Docker Compose \citep{docker_2026} y se destinan a ejecutarse como
contenedores Podman \citep{podman_2026} en el servidor que la Universidad del
Bío-Bío dispone para proyectos de título. Esta arquitectura responde al
requerimiento del profesor guía de contar con un backend y una base de datos
propios, y entrega control total sobre los datos y su cifrado a cambio de
asumir la operación del servidor. La Tabla~\ref{tab:herramientas} detalla el
propósito de cada herramienta.

\begin{table}[H]
\centering
\small
\tablazebra
\caption{Tecnologías y herramientas utilizadas en el desarrollo.}
\label{tab:herramientas}
\begin{tabularx}{\textwidth}{P{0.75}P{1.25}}
\toprule
\filaencabezado \encab{Tecnología o herramienta} & \encab{Propósito en el proyecto} \\
\midrule
Flutter y Dart & Construcción de la aplicación móvil para Android e iOS. \\
Riverpod & Gestión del estado y separación entre la lógica de datos y la
presentación \citep{riverpod_2026}. \\
go\_router & Navegación entre pantallas y redirección según el estado de la
sesión \citep{go_router_2026}. \\
Node.js 22 y Express 5 & API REST que atiende a la aplicación, valida las
solicitudes y custodia la clave de cifrado de cada usuario. \\
PostgreSQL 15 & Base de datos relacional que almacena los datos del cuidador, los
pacientes y sus registros como texto cifrado. \\
jsonwebtoken y bcryptjs & Sesión con tokens JWT y almacenamiento de las
contraseñas como \emph{hash} bcrypt. \\
nodemailer y SMTP & Envío de los correos de verificación y de recuperación de
acceso. \\
ffmpeg & Generación de las miniaturas de los videos de la biblioteca. \\
Docker y Docker Compose & Empaquetado y ejecución local de la API y la base de
datos. \\
Podman & Ejecución de los contenedores en el servidor de la universidad (destino
del despliegue). \\
DBeaver & Inspección de la base de datos durante el desarrollo. \\
cryptography y flutter\_secure\_storage & Cifrado AES-256-GCM de los datos
sensibles y resguardo de la clave en el dispositivo \citep{cryptography_dart_2026}. \\
flutter\_local\_notifications & Programación de los recordatorios como
notificaciones locales. \\
pdf y excel & Generación de los informes exportables del historial. \\
Git y GitHub & Control de versiones y respaldo del código, con una rama por
semana de trabajo. \\
Figma & Diseño de los prototipos de interfaz. \\
flutter\_test & Ejecución de la suite de pruebas automatizadas de la aplicación. \\
Jest y supertest & Pruebas de la API contra una base PostgreSQL real. \\
\bottomrule
\end{tabularx}
\end{table}

\section[Definiciones, Siglas y Abreviaciones]{Definiciones, siglas y abreviaciones}

\begingroup
\small
\tablazebra
\begin{xltabular}{\textwidth}{P{0.50}P{1.50}}
    \caption{Definiciones, siglas y abreviaciones.}
    \label{tab:siglas}\\
    \toprule
    \filaencabezado \encab{Término} & \encab{Definición} \\
    \midrule
    \endfirsthead
    \rowcolor{white}\multicolumn{2}{@{}l}{\footnotesize\itshape Tabla \thetable{} (continuación)} \\
    \toprule
    \filaencabezado \encab{Término} & \encab{Definición} \\
    \midrule
    \endhead
    \midrule
    \rowcolor{white}\multicolumn{2}{r@{}}{\footnotesize\itshape Continúa en la página siguiente} \\
    \endfoot
    \bottomrule
    \endlastfoot
    AES-256-GCM & Algoritmo de cifrado simétrico con clave de 256 bits en modo
    Galois/Counter, que además de cifrar verifica la integridad del dato. \\
    APK & Archivo de instalación de aplicaciones para Android. \\
    Backlog & Lista priorizada de historias de usuario pendientes de desarrollo. \\
    BPMN & Business Process Model and Notation, notación para modelar procesos
    de negocio. \\
    Cuidador principal & Persona que asume el cuidado diario del paciente en el
    hogar y usuario de la aplicación. \\
    CU & Caso de uso. \\
    ESAS & Edmonton Symptom Assessment System, escala que valora la intensidad de
    los síntomas de 0 a 10. \\
    HMAC & Hash-based Message Authentication Code, código de autenticación de
    mensajes basado en una función \emph{hash}; el servidor lo usa para guardar
    el correo de respaldo sin almacenarlo en claro. \\
    HTTP / HTTPS & Protocolo de transferencia de hipertexto, sin cifrar o con
    cifrado de transporte, que usa la aplicación para comunicarse con la API. \\
    HU & Historia de usuario. \\
    JSON & JavaScript Object Notation, formato de texto para intercambiar datos
    entre la aplicación y la API. \\
    JWT & JSON Web Token, token firmado que identifica la sesión del usuario en
    cada solicitud a la API. \\
    Motor de reglas & Componente que aplica condiciones predefinidas para
    clasificar un registro o elegir una respuesta, sin inteligencia artificial
    generativa. \\
    Offline-first & Enfoque de diseño en que la aplicación funciona sin conexión
    y sincroniza sus datos cuando recupera la red. \\
    Podman & Motor de contenedores sin demonio, compatible con Docker, que se
    usa en el servidor de la universidad. \\
    PINDA & Programa Infantil Nacional de Drogas Antineoplásicas del Ministerio
    de Salud de Chile. \\
    REST / API & Estilo de interfaz de servicios web sobre HTTP / interfaz de
    programación de aplicaciones. \\
    RF / RNF & Requerimiento funcional / requerimiento no funcional. \\
    Semaforización & Clasificación del nivel de riesgo en tres colores: verde
    (normal), amarillo (alerta) y rojo (crítico). \\
    SMTP & Simple Mail Transfer Protocol, protocolo con que el servidor envía los
    correos de verificación y recuperación. \\
    Sprint & Periodo fijo de trabajo al final del cual se entrega una versión
    funcional del producto. \\
    TRL & Technology Readiness Level. El nivel~3 corresponde a una prueba de
    concepto validada en un entorno de desarrollo. \\
    UCP & Use Case Points, técnica de estimación de tamaño basada en casos de
    uso. \\
    VAN / TIR & Valor actual neto / tasa interna de retorno. \\
\end{xltabular}
\endgroup

\addtocontents{toc}{\protect\setcounter{tocdepth}{2}}
